\section{Completion of the CM proof}\label{sec:assembly}

The theta construction gives a nonzero period on a compact ball
quotient.  The finite-family extraction corollary expresses the one-form
classes in that period through rational maps from the prescribed CM
modules, after passage to sufficiently deep levels.  We put these
statements on one projective surface and apply the criterion of
Section~\ref{sec:surface}.

\begin{proof}[Proof of Theorem~\ref{thm:switch}]
By Lemma~\ref{lem:switch-normalization}, it suffices to treat the case
\(s\geq2\) and \(v_i(1)=1\) for all four functions.  Use the common
Galois CM field \(E\), the active embedding \(\mu=1\), and the line
signs
\[
 d_i(\tau)=v_i(\tau^{-1}).
\]
They satisfy \(d_1+d_2=d_3+d_4\) and \(d_i(\mu)=1\).  Choose the
rank-three Hermitian space \(V\), the ball \(D\), and the rational
orthogonal decomposition \(W=W_1\perp W_2\) prescribed by \(d_1,d_2\),
as in Section~\ref{sec:theta}.

Proposition~\ref{prop:theta-top}(ii) supplies a total finite theta
input whose form for this decomposition is nonzero.  With that input
fixed, Corollary~\ref{cor:theta-switch-continuity} supplies a second
rational orthogonal decomposition and a nonzero mixed period.
Write \(X'\) for the individual-factor cover
constructed there, and \(\xi_1,\xi_2\) for the two pulled-back endpoint
forms.  On this common compact principal quotient they have finite expansions
\[
 \xi_1=\sum_{a\in I}c_a\,
       \Theta_1(\phi_{1a})\wedge\Theta_2(\phi_{2a}),
 \qquad
 \xi_2=\sum_{b\in J}c'_b\,
       \Theta_3(\phi_{3b})\wedge\Theta_4(\phi_{4b}),
\]
where the constants account for the chosen constant-line
trivializations, and
\[
 \int_{X'}\xi_1\wedge\overline{\xi_2}\ne0.
\]
All four families of individual inputs here are finite.  The quotient
and these expansions may be pulled to any deeper principal level;
the integral is then multiplied by the degree of the cover.

Apply Corollary~\ref{cor:finite-family-extraction} to these four families,
with \(X_{\mathrm{base}}=X'\) and the fixed \(E,V,D,\Lambda\).
It gives a common smooth connected projective principal cover
\(\pi:Y\to X'\) on which all four \(U(v_i)_1\)-span conclusions hold;
the auxiliary primes may differ.

Denote the pulled-back one-form classes by
\[
 \vartheta_{ia}=[\pi^*\Theta_i(\phi_{ia})]\in H^1(Y,\C),
\]
with \(a\in I\) for \(i=1,2\) and \(a\in J\) for \(i=3,4\).
These classes are represented by holomorphic one-forms.  Their cup
products compute the integrals of the wedge products, and the
covering formula gives
\begin{equation}\label{eq:assembly-period}
\begin{split}
 0&\ne\int_Y\pi^*\xi_1\wedge\overline{\pi^*\xi_2}\\
  &=\sum_{a\in I,\ b\in J}c_a\overline{c'_b}\,
    \left\langle
      \vartheta_{1a}\cup\vartheta_{2a}\cup
      \overline{\vartheta_{3b}}\cup\overline{\vartheta_{4b}},[Y]
    \right\rangle .
\end{split}
\end{equation}
Choose \(a,b\) for which the displayed cup-product period is nonzero,
and write its four holomorphic factors as \(\vartheta_i\).

Corollary~\ref{cor:finite-family-extraction} gives
\[
 \vartheta_i\in\operatorname{span}_{\C}
 \{h(U(v_i)_1):
       h\in\Hom_{\mathrm{Hdg}}(U(v_i),H^1(Y,\Q))\},
 \qquad 1\leq i\leq4.
\]
The maps are rational, so they commute with scalar complex
conjugation, which sends \(U(v_i)_1\) to \(U(v_i)_c\).  Consequently
\[
 \overline{\vartheta_i}\in\operatorname{span}_{\C}
 \{h(U(v_i)_c):
       h\in\Hom_{\mathrm{Hdg}}(U(v_i),H^1(Y,\Q))\},
 \qquad i=3,4.
\]
The four classes
\(\vartheta_1,\vartheta_2,\overline{\vartheta_3},
\overline{\vartheta_4}\) thus satisfy
Corollary~\ref{cor:surface-span} on the same projective surface \(Y\),
with the source labels \((1,1,c,c)\) and the nonzero period selected
from Equation~\eqref{eq:assembly-period}.  The normalized switch line
is therefore algebraic.  Lemma~\ref{lem:switch-normalization} returns
every switch to this case or
to a product of algebraic two-factor classes, and proves the theorem.
\end{proof}

\begin{proof}[Proof of Theorem~\ref{thm:cm}]
The four-factor switch theorem is now proved for every Galois CM field.
Proposition~\ref{prop:switch-reduction} consequently places every
rational Hodge class on a CM abelian variety in its rational cycle
span.  Its final scalar descent takes place on the original variety.
Finite products and powers are CM, so the same conclusion covers them.
\end{proof}
